% !TeX spellcheck = ru_RU
% !TEX root = vkr.tex

\section*{Введение}

Многие прикладные алгоритмы~--- например, в областях искусственного интеллекта и анализа графов~--- естественным образом выражаются в терминах линейной алгебры.
Такой подход позволяет использовать готовые библиотеки, например, \textsc{BLAS}~\cite{blas} для плотных матриц и \textsc{GraphBLAS}~\cite{graphblas} для разреженных, использующие в том числе большие возможности линейной алгебры по распараллеливанию вычислений.

Если в случае плотных матриц основными ограничениями являются объём памяти, необходимый для хранения, и скорость передачи данных, то в случае разреженных основными проблемами являются нелокальные обращения к памяти и непредсказуемое количество независимых подзадач, в силу принципов проектирования современных высокопроизводительных CPU и GPU.

\INs{}~--- модель вычислений, которая была описана Yves Lafont в 1989 году~\cite{lafontInteractionNets1989}.
В этой модели программа представляется в виде графа, а одним из важных свойств является естественная поддержка нерегулярного параллелизма.

Существующие реализации исполнителей \INs{} демонстрируют скорость исполнения лишь на базовых конструкциях (сортировка, обход списков).
В рамках проекта \Lamagraph{}\footnote{Репозитории проекта доступны по ссылке: \url{https://github.com/Lamagraph/} (дата обращения: \DTMdate{2026-05-18})} изучаются возможности по реализации алгоритмов разреженной линейной алгебры на \INs{}.
Для оценки применимости \INs{} был разработан небольшой инструмент для сравнения производительности различных реализаций алгоритмов линейной алгебры.

В данной работе предлагается произвести рефакторинг данного инструмента, путём унификации общего кода, выделения единообразного интерфейса, а также создания документации по расширению другими реализациями.
